\clearpage\SourceLeaf{190}
\begin{center}
{\Large\textls{REVIEW.}}
\end{center}

\bigskip
\noindent\hrulefill

\begin{center}
Göttingische gelehrte Anzeigen. 1822 October 28.
\end{center}

\noindent\hrulefill
\bigskip

\begin{center}
\textls{Marburg.}
\end{center}

\emph{Theorie der Parallelen}, by \textsc{Carl Reinhard Müller}, \emph{Doctor of Philosophy, extraordinary professor of mathematics} and\,so\,on. 1822. 40 pages in quarto.

Six years ago the reviewer expressed in these pages his conviction that all previous attempts to prove the theory of parallel lines rigorously or to fill the gap in Euclidean geometry had brought us no closer to this goal. He can only extend this judgement to all later attempts that have come to his attention. Even so, some of these attempts remain worth reading for those interested in geometry because of the ingenuity they show, even though they miss their chief aim. In this respect the reviewer believes that the present short work, issued on the occasion of a school examination, deserves special notice. A detailed account of the author's ingenious reasoning would take up too much space in these pages and would be superfluous, since the work itself deserves to be read. Yet it has a weak point, as do all the other attempts, and this review is meant to bring it out.

We find this weak point on p.\,15 in the proof of the theorem of Article 15. This theorem is the very nerve of the whole theory, which falls as soon as it cannot be rigorously proved. We therefore state the theorem first. Every reader can readily draw the corresponding figure.

If each angle at the base $ON$ of an isosceles triangle is greater than the angle at the vertex $A$, and if at $O$ one constructs on the side $OA$
